\chapter{Detection measures: ROC, pAUC, TPR at a fixed FPR, objects caught and chance}\label{ch:measures}

\begin{onething}
When anomalies are 0.1--1\% of a scene's pixels, the area under the ROC curve is decided almost entirely by how
ordinary background pixels are ranked against each other far down the list, where no analyst ever looks; so the
repository also measures the \emph{top} of the ranking - partial AUC up to 1\% false alarms, the detection rate at
0.1\% false alarms, and how many objects land in a fixed review budget - and reads every one of them against what
random scores would get.
\end{onething}

\section{Why a new set of measures was needed}\label{sec:measures:why}

Every chapter so far has ended with a \emph{score map}: one number per pixel, high where the pixel looks
anomalous (\cref{ch:residualmaha}). To compare two detectors you need a way to turn a score map and a
\index{ground truth}\emph{ground-truth mask} - a map marking which pixels a human labelled as anomalous - into a
number. Until 2026-10-16 the repository used one number for this, the AUC returned by
\code{compute_roc} in \file{app/utils/anomaly_detection/scoring.py}. On that date a second family of measures
arrived beside it. Since 2026-10-17 they live in \file{app/evaluation/metrics.py}, keyed \code{pauc_0.01},
\code{tpr_at_0.001}, \code{hit_1} and \code{hit_0.1}, and are tested by hand-worked numbers in
\file{tests/test_evaluation/test_metrics.py}.

The problem the new measures solve is easy to state and surprisingly easy to miss. A PRISMA scene has about a
million valid pixels. Suppose a thousand of them are targets\tagA{} (0.1\%, the order of the labelled bench
scenes). An analyst can review perhaps the top 1\% of the ranking: ten thousand pixels. Anything the detector does
to the other 990,000 pixels is invisible to that analyst. A good measure must therefore be sensitive to the top of
the list and indifferent to the bottom. The AUC is the opposite, as \cref{sec:measures:bulk} proves with numbers.

\begin{analogy}
A search engine returns a million results for your query. You read the first page. A ranking metric that gives the
same weight to the order of results 500,000 and 500,001 as to results 1 and 2 will happily prefer an engine whose
first page is junk, as long as it sorts the junk at the back slightly better. The analyst is the reader of the
first page; the measures of this chapter are page-one measures.
\end{analogy}

You need two ideas from earlier: that anomalies are rare (\cref{ch:problem}) and the idea of a threshold on a
score (\cref{ch:residualmaha}). Nothing else.

\section{Thresholds and the confusion matrix}\label{sec:measures:confusion}

\begin{definition}[Positive, negative, flagged]
A \index{positive}\emph{positive} (or \emph{target}) is a valid pixel the ground truth marks as anomalous. A
\index{negative}\emph{negative} (or \emph{background} pixel) is every other valid pixel. Given a
\index{threshold}\emph{threshold} $\tau$, a pixel is \emph{flagged} when its score $s \ge \tau$.
\end{definition}

Every pixel then falls into exactly one of four boxes, and the table of their counts is called the
\index{confusion matrix}\emph{confusion matrix}:

\begin{center}
\begin{tabular}{lcc}
\toprule
 & truth: target & truth: background \\
\midrule
flagged ($s \ge \tau$) & true positive (TP) & false positive (FP) \\
not flagged ($s < \tau$) & false negative (FN) & true negative (TN) \\
\bottomrule
\end{tabular}
\end{center}

Write $P$ for the number of targets and $N$ for the number of background pixels. Each column adds up to its total:
$\mathrm{TP} + \mathrm{FN} = P$ and $\mathrm{FP} + \mathrm{TN} = N$. Two rates summarise the table.

\begin{definition}[TPR and FPR]
The \index{true-positive rate}\emph{true-positive rate} (TPR, also the detection rate) and the
\index{false-positive rate}\emph{false-positive rate} (FPR, also the false-alarm rate) at threshold $\tau$ are
\begin{equation}\label{eq:measures:rates}
\mathrm{TPR}(\tau) = \frac{\mathrm{TP}(\tau)}{P}, \qquad \mathrm{FPR}(\tau) = \frac{\mathrm{FP}(\tau)}{N}.
\end{equation}
\end{definition}

Both rates are divided by their own column total, which is the point: FPR does not care how many targets there
are, and TPR does not care how much background there is. Both are \emph{non-increasing} functions of $\tau$: raise
the threshold and you can only flag fewer pixels, never more.

\section{Six pixels by hand: the ROC curve}\label{sec:measures:six}

\begin{definition}[ROC curve]
The \index{ROC curve}\emph{receiver operating characteristic} (ROC) curve is the set of points
$(\mathrm{FPR}(\tau), \mathrm{TPR}(\tau))$ as $\tau$ slides from above the highest score (nothing flagged, the
point $(0,0)$) down to the lowest score (everything flagged, the point $(1,1)$), joined by straight segments.
\end{definition}

You only need to try thresholds equal to the scores themselves: between two neighbouring scores nothing changes.

\begin{worked}[label={we:measures:six}]{The ROC of six pixels}
Do this one on paper. Six valid pixels, three targets (T) and three background pixels (B), so $P = 3$, $N = 3$:
\begin{center}
\begin{tabular}{lcccccc}
\toprule
pixel & 1 & 2 & 3 & 4 & 5 & 6 \\
score & 0.95 & 0.85 & 0.70 & 0.60 & 0.40 & 0.20 \\
truth & T & B & T & T & B & B \\
\bottomrule
\end{tabular}
\end{center}
\textbf{Step 1 - sort by score, highest first.} They are already sorted.

\textbf{Step 2 - sweep the threshold down one pixel at a time}, counting what has been flagged:
\begin{center}
\begin{tabular}{lcccc}
\toprule
threshold $\tau$ & TP & FP & $\mathrm{TPR} = \mathrm{TP}/3$ & $\mathrm{FPR} = \mathrm{FP}/3$ \\
\midrule
above 0.95 & 0 & 0 & 0 & 0 \\
0.95 (T) & 1 & 0 & 1/3 & 0 \\
0.85 (B) & 1 & 1 & 1/3 & 1/3 \\
0.70 (T) & 2 & 1 & 2/3 & 1/3 \\
0.60 (T) & 3 & 1 & 1 & 1/3 \\
0.40 (B) & 3 & 2 & 1 & 2/3 \\
0.20 (B) & 3 & 3 & 1 & 1 \\
\bottomrule
\end{tabular}
\end{center}
\textbf{Step 3 - list the ROC points} $(\mathrm{FPR}, \mathrm{TPR})$:
\[
(0,0),\ (0,\tfrac13),\ (\tfrac13,\tfrac13),\ (\tfrac13,\tfrac23),\ (\tfrac13,1),\ (\tfrac23,1),\ (1,1).
\]
A target moves the curve \emph{up} by $1/P = 1/3$; a background pixel moves it \emph{right} by $1/N = 1/3$. The
curve is a staircase. \Cref{fig:measures:sixroc} draws it.
\end{worked}

The rule in the last sentence is general and worth stating on its own.

\begin{keypoint}
Walking down the sorted list, every target is a step up of $1/P$ and every background pixel a step right of $1/N$;
the ROC curve is the path that walk traces.
\end{keypoint}

\section{The area under the curve, by the trapezoid rule}\label{sec:measures:auc}

\begin{definition}[AUC]
The \index{AUC}\emph{area under the ROC curve} (AUC, or ROC-AUC) is the area between the curve and the FPR axis,
from FPR 0 to FPR 1. A perfect detector, which ranks every target above every background pixel, goes straight up to
$(0,1)$ and then right: AUC 1. A detector that ranks at random follows the diagonal $\mathrm{TPR} = \mathrm{FPR}$:
AUC $1/2$.
\end{definition}

Because the curve is made of straight segments, its area is a sum of trapezoids. A \index{trapezoid rule}\emph{trapezoid}
between points $(x_0, y_0)$ and $(x_1, y_1)$ has area width times average height:
\begin{equation}\label{eq:measures:trap}
\text{area} = (x_1 - x_0)\,\frac{y_0 + y_1}{2}.
\end{equation}
This is exactly what \code{auc_exact} in \file{app/evaluation/metrics.py} computes, one line:
\begin{lstlisting}[basicstyle=\ttfamily\footnotesize]
def auc_exact(fpr, tpr):     # fpr, tpr: (K,) ROC points from exact_roc, starting at (0, 0)
    return float(np.sum((fpr[1:] - fpr[:-1]) * (tpr[1:] + tpr[:-1]) / 2))
\end{lstlisting}

\begin{worked}[label={we:measures:sixauc}]{The AUC of the six pixels, strip by strip}
Apply \cref{eq:measures:trap} to each consecutive pair of points of \cref{we:measures:six}:
\begin{center}
\begin{tabular}{lccc}
\toprule
segment & width & heights & area \\
\midrule
$(0,0) \to (0,\frac13)$ & $0$ & $0, \frac13$ & $0$ \\
$(0,\frac13) \to (\frac13,\frac13)$ & $\frac13$ & $\frac13, \frac13$ & $\frac13 \times \frac13 = \frac19$ \\
$(\frac13,\frac13) \to (\frac13,\frac23)$ & $0$ & & $0$ \\
$(\frac13,\frac23) \to (\frac13,1)$ & $0$ & & $0$ \\
$(\frac13,1) \to (\frac23,1)$ & $\frac13$ & $1, 1$ & $\frac13$ \\
$(\frac23,1) \to (1,1)$ & $\frac13$ & $1, 1$ & $\frac13$ \\
\bottomrule
\end{tabular}
\end{center}
Add the column:
\begin{align*}
\text{AUC} &= \tfrac19 + \tfrac13 + \tfrac13 \\
&= \tfrac19 + \tfrac39 + \tfrac39 \\
&= \tfrac79 = 0.7778.
\end{align*}
The repository's \code{auc_exact} on these six pixels returns $0.7778$\tagC{} (the book's check script
\file{research/textbook/tools/compute_measures.py}).
\end{worked}

\begin{figure}[htbp]
\centering
\begin{tikzpicture}
\begin{axis}[width=8.2cm,height=8.2cm,xmin=0,xmax=1,ymin=0,ymax=1.02,
  xlabel={false-positive rate (FPR)},ylabel={true-positive rate (TPR)},
  xtick={0,0.3333,0.5,0.6667,1},xticklabels={0,$\frac13$,$\frac12$,$\frac23$,1},
  ytick={0,0.3333,0.6667,1},yticklabels={0,$\frac13$,$\frac23$,1},grid=major,grid style={gray!20}]
\addplot[fill=loss!30,draw=none] coordinates {(0,0) (0,0.3333) (0.3333,0.3333) (0.3333,1) (0.5,1) (0.5,0)} \closedcycle;
\addplot[dashed,maskgrey,thick] coordinates {(0,0) (1,1)};
\addplot[spat,very thick,mark=*,mark options={fill=white}] coordinates {(0,0) (0,0.3333) (0.3333,0.3333) (0.3333,0.6667) (0.3333,1) (0.6667,1) (1,1)};
\draw[loss,thick,dashed] (axis cs:0.5,0) -- (axis cs:0.5,1);
\node[anchor=west,font=\small] at (axis cs:0.52,0.22) {chance (diagonal)};
\node[anchor=west,font=\small,text=loss!70!black] at (axis cs:0.03,0.12) {pAUC region};
\node[anchor=west,font=\small,text=loss!70!black] at (axis cs:0.03,0.05) {$\mathrm{FPR}\le 0.5$};
\node[anchor=south west,font=\small,text=spat] at (axis cs:0.34,0.62) {after 0.70 (T)};
\node[anchor=north west,font=\small,text=spat] at (axis cs:0.01,0.32) {after 0.95 (T)};
\end{axis}
\end{tikzpicture}
\caption{The ROC of the six pixels of \cref{we:measures:six} (indigo staircase), the chance diagonal (grey,
dashed), and the region counted by a partial AUC with limit $\mathrm{FPR} \le 0.5$ (amber). Drawn to scale from
the worked numbers. The full area under the staircase is $7/9$; the amber area is $5/18$, which normalised by the
limit gives $5/9$ (\cref{we:measures:pauc}).}
\label{fig:measures:sixroc}
\end{figure}

\section{AUC is a probability: the Mann--Whitney count}\label{sec:measures:mw}

There is a second way to read the same number, and it explains everything that follows.

\begin{definition}[Pair reading of AUC]
Pick one target at random and one background pixel at random. The \index{Mann--Whitney statistic}\emph{pair
reading} of AUC is the probability that the target scores higher, counting a tie as half:
\begin{equation}\label{eq:measures:mw}
\mathrm{AUC} = \frac{1}{PN}\sum_{\text{targets } t}\ \sum_{\text{background } b}
\Big(\mathbf{1}[s_t > s_b] + \tfrac12\,\mathbf{1}[s_t = s_b]\Big).
\end{equation}
Here $\mathbf{1}[\cdot]$ is the \emph{indicator}: 1 when the statement in brackets is true, 0 when false. The double
sum counts winning pairs; it is the \emph{Mann--Whitney U count}, and dividing by the $PN$ pairs turns it into a
probability \citep{hanley1982meaning}.
\end{definition}

\begin{worked}[label={we:measures:mw}]{The six pixels, counted as pairs}
There are $P \times N = 3 \times 3 = 9$ (target, background) pairs. For each target, count the background pixels it
beats:
\begin{align*}
\text{target } 0.95 &\text{ vs } 0.85, 0.40, 0.20: \text{ beats all three} &&\to 3\\
\text{target } 0.70 &\text{ vs } 0.85, 0.40, 0.20: \text{ loses to } 0.85 &&\to 2\\
\text{target } 0.60 &\text{ vs } 0.85, 0.40, 0.20: \text{ loses to } 0.85 &&\to 2\\
\text{wins} &= 3 + 2 + 2 = 7, &&\text{AUC} = 7/9 = 0.7778.
\end{align*}
The same $7/9$ as the trapezoids of \cref{we:measures:sixauc}.
\end{worked}

That the two readings agree is not a coincidence. Here is the proof, in words first and then in symbols.

\paragraph{Proof that the area equals the pair count (no ties).} By the key point of \cref{sec:measures:six}, the
curve moves right only at background pixels, by $1/N$ each, and moves up only at targets. So the area under the
staircase is a sum of rectangles, one per background pixel $b$: width $1/N$, and height equal to the TPR at the
moment $b$ is reached in the walk. That TPR is the share of targets already passed, which is the share of targets
scoring above $b$:
\begin{align*}
\text{area} &= \sum_{b} \frac1N \cdot \mathrm{TPR}\text{ when $b$ is reached}
  && \text{one rectangle per background pixel}\\
&= \sum_{b} \frac1N \cdot \frac{\#\{t : s_t > s_b\}}{P}
  && \text{TPR = share of targets above $b$}\\
&= \frac{1}{PN} \sum_b \sum_t \mathbf{1}[s_t > s_b]
  && \text{write the count as a sum of indicators}
\end{align*}
which is \cref{eq:measures:mw} with no ties. Check it on the six pixels: the background pixels are reached at
heights $\frac13$ (pixel 0.85), $1$ (pixel 0.40) and $1$ (pixel 0.20), so the area is
$\frac13(\frac13 + 1 + 1) = \frac13 \cdot \frac73 = \frac79$.

\paragraph{With ties.} When a target and a background pixel share a score, the threshold flags both at once, and
the curve moves diagonally: right by $1/N$ and up by $1/P$ in one segment. A trapezoid under that diagonal step
gets the average of the heights before and after - exactly half the target's contribution, the $\frac12$ in
\cref{eq:measures:mw}.

\begin{worked}[label={we:measures:tie}]{A tie, both ways}
Four pixels: $0.9$ (T), $0.5$ (T), $0.5$ (B), $0.1$ (B), so $P = N = 2$. The function \code{exact_roc} puts one ROC
point at the end of each run of tied scores (\cref{sec:measures:exact}), so the two 0.5s are flagged together.
\begin{align*}
&\text{points: } (0,0),\ (0, \tfrac12),\ (\tfrac12, 1),\ (1, 1) \quad \text{(the middle segment is diagonal)}\\
&\text{trapezoids: } 0 + \tfrac12 \cdot \tfrac{\frac12 + 1}{2} + \tfrac12 \cdot \tfrac{1 + 1}{2}\\
&\qquad = 0 + \tfrac12 \cdot \tfrac34 + \tfrac12 = \tfrac38 + \tfrac48 = \tfrac78 = 0.875\\
&\text{pairs: } 0.9 \text{ beats } 0.5, 0.1 \to 2; \quad 0.5 \text{ ties } 0.5 \to \tfrac12\\
&\qquad \text{and beats } 0.1 \to 1; \quad \text{total } 2 + 1.5 = 3.5\\
&\text{AUC} = 3.5 / 4 = 0.875.
\end{align*}
\end{worked}

\section{Why AUC is decided by the bulk}\label{sec:measures:bulk}

Read \cref{eq:measures:mw} again with the real scene in mind: $P = 1{,}000$ targets, $N = 999{,}000$ background
pixels\tagA{}. AUC averages over every one of the $999{,}000{,}000$ pairs. A target ranked 500,000th still beats
the 499,000 or so background pixels below it, and every one of those wins counts. \textbf{Almost all of the AUC is
earned by putting targets above pixels nobody will ever look at.} The analyst cares about the region
$\mathrm{FPR} \in [0, 0.01]$, which is 1\% of the curve's width.

Two invented detectors (from the course, part 16) make it concrete. Their ROC curves, as straight segments:
\begin{align*}
\text{A}: &\ (0,0) \to (0.01, 0.40) \to (0.30, 0.80) \to (1,1)\\
\text{B}: &\ (0,0) \to (0.01, 0.02) \to (0.10, 0.80) \to (1,1)
\end{align*}

\begin{worked}[label={we:measures:ab}]{Two detectors whose AUC and top-of-list disagree}
Trapezoids for A:
\begin{align*}
0.01 \times \tfrac{0 + 0.40}{2} &= 0.01 \times 0.20 = 0.002\\
0.29 \times \tfrac{0.40 + 0.80}{2} &= 0.29 \times 0.60 = 0.174\\
0.70 \times \tfrac{0.80 + 1.00}{2} &= 0.70 \times 0.90 = 0.630\\
\text{AUC}_A &= 0.002 + 0.174 + 0.630 = 0.806.
\end{align*}
Trapezoids for B:
\begin{align*}
0.01 \times \tfrac{0 + 0.02}{2} &= 0.01 \times 0.01 = 0.0001\\
0.09 \times \tfrac{0.02 + 0.80}{2} &= 0.09 \times 0.41 = 0.0369\\
0.90 \times \tfrac{0.80 + 1.00}{2} &= 0.90 \times 0.90 = 0.810\\
\text{AUC}_B &= 0.0001 + 0.0369 + 0.810 = 0.847.
\end{align*}
B has the higher AUC. But at 1\% false alarms A has found 40\% of the targets and B 2\%.
\end{worked}

\coursefig[0.95\linewidth]{ad-16-roc-top}{The course's figure for this chapter. (a) The five-pixel ROC of the
course's own example (targets at 0.9 and 0.7 among five pixels), AUC $5/6$. (b) Detectors A and B: B has the larger
total area; inside the first 1\% of false alarms (the amber strip, stretched, not to scale) A is twenty times
higher.}{fig:measures:rocTop}

\begin{figure}[htbp]
\centering
\begin{tikzpicture}
\begin{groupplot}[group style={group size=2 by 1,horizontal sep=1.6cm},width=6.4cm,height=6cm,
  ymin=0,ymax=1.02,xlabel={FPR},grid=major,grid style={gray!20}]
\nextgroupplot[xmin=0,xmax=1,ylabel={TPR},title={\small whole curve}]
\addplot[dashed,maskgrey] coordinates {(0,0) (1,1)};
\addplot[spec,very thick,mark=*,mark size=1.5pt] coordinates {(0,0) (0.01,0.40) (0.30,0.80) (1,1)};
\addplot[spat,very thick,dashed,mark=square*,mark size=1.5pt] coordinates {(0,0) (0.01,0.02) (0.10,0.80) (1,1)};
\node[font=\small,text=spec,anchor=west] at (axis cs:0.42,0.30) {A: AUC 0.806};
\node[font=\small,text=spat,anchor=west] at (axis cs:0.42,0.20) {B: AUC 0.847};
\nextgroupplot[ymax=0.5,xmin=0,xmax=0.012,xtick={0,0.005,0.01},xticklabels={0,0.005,0.01},title={\small first 1\% only, to scale},
  scaled x ticks=false]
\addplot[fill=loss!25,draw=none] coordinates {(0,0) (0.01,0.40) (0.01,0)} \closedcycle;
\addplot[dashed,maskgrey] coordinates {(0,0) (0.012,0.012)};
\addplot[spec,very thick] coordinates {(0,0) (0.01,0.40) (0.012,0.4276)};
\addplot[spat,very thick,dashed] coordinates {(0,0) (0.01,0.02) (0.012,0.0373)};
\node[font=\small,text=spec,anchor=west] at (axis cs:0.0003,0.46) {A: pAUC 0.20};
\node[font=\small,text=spat,anchor=west] at (axis cs:0.0003,0.41) {B: pAUC 0.01};
\end{groupplot}
\end{tikzpicture}
\caption{Detectors A (teal) and B (indigo, dashed) of \cref{we:measures:ab}, drawn to scale. Left: the whole curve,
where B's larger area comes from FPR 0.01--0.30. Right: the first 1.2\% of false alarms, where A's area (amber, up to
FPR 0.01) is twenty times B's. The chance diagonal (grey) is nearly flat on the right-hand scale.}
\label{fig:measures:ab}
\end{figure}

\begin{keypoint}
A ranking measure averaged over the whole list can prefer the detector that is useless where it is used.
\end{keypoint}

\section{Partial AUC up to a false-alarm limit}\label{sec:measures:pauc}

\begin{definition}[Normalised partial AUC]
The \index{partial AUC}\emph{partial AUC} up to an FPR limit $m$ is the area under the ROC between FPR 0 and FPR
$m$. The repository divides it by $m$ so that a perfect detector scores 1:
\begin{equation}\label{eq:measures:pauc}
\mathrm{pAUC}_m = \frac{1}{m}\int_0^{m} \mathrm{TPR}(f)\,df.
\end{equation}
The repository's headline is $m = 0.01$, key \code{pauc_0.01}, configured in an evaluation config as
\texttt{\{"kind": "pauc", "max\_fpr": 0.01\}}.
\end{definition}

\begin{sourcenote}[: what pAUC is normalised by]
\file{app/evaluation/metrics.py}, module docstring and function \code{pauc} (lines 31--36): ``area under the EXACT
ROC over FPR 0..max\_fpr, divided by max\_fpr (1 perfect, max\_fpr/2 random)''. The default \code{max_fpr} is
0.01. The normalisation is a plain division by the width $m$. It is \emph{not} the McClish standardisation used by
some libraries (scikit-learn's \code{roc_auc_score(max_fpr=...)} applies McClish); derived below so you can tell the
two apart when you compare with outside numbers.
\end{sourcenote}

\paragraph{Its chance value, derived.} A random detector's ROC is the diagonal $\mathrm{TPR}(f) = f$. Put that
into \cref{eq:measures:pauc}:
\begin{align*}
\mathrm{pAUC}_m^{\text{chance}} &= \frac1m \int_0^m f\,df && \text{TPR equals FPR on the diagonal}\\
&= \frac1m \cdot \Big[\tfrac12 f^2\Big]_0^m && \text{the integral of $f$ is $\tfrac12 f^2$}\\
&= \frac1m \cdot \tfrac12 m^2 && \text{evaluate at $m$, minus at $0$}\\
&= \frac{m}{2}. && \text{cancel one $m$}
\end{align*}
For $m = 0.01$: chance is $0.01/2 = 0.005$. Geometrically it is a triangle of base 0.01 and height 0.01, area
$0.01 \times 0.01 / 2 = 0.00005$, divided by $0.01$. \textbf{Random pAUC is 0.005, not 0.5.} The test
\code{test_perfect_random_and_sklearn} checks that random scores stay below 0.05; on two million random pixels the
book's script measured $0.0052$, $0.0055$ and $0.0057$ in three draws\tagC{}.

\paragraph{The McClish standardisation, for comparison.} \citet{mcclish1989analyzing} rescales the raw partial
area $A$ (not divided by $m$) so that chance maps to $\frac12$ and perfect to 1, like a full AUC. The smallest
sensible raw area is the chance triangle $A_{\min} = m^2/2$; the largest is the full rectangle $A_{\max} = m$:
\begin{equation}\label{eq:measures:mcclish}
\mathrm{pAUC}^{\text{McClish}} = \frac12\left(1 + \frac{A - A_{\min}}{A_{\max} - A_{\min}}\right).
\end{equation}
Check the ends: $A = A_{\min}$ gives $\frac12(1 + 0) = \frac12$, and $A = A_{\max}$ gives $\frac12(1 + 1) = 1$.
For a repository pAUC of 0.170 at $m = 0.01$ (FT04 with cloud excluded, \cref{ch:bootstrap}):
\begin{align*}
A &= 0.170 \times 0.01 = 0.0017 && \text{undo the division by $m$}\\
A_{\min} &= 0.01^2 / 2 = 0.00005, \quad A_{\max} = 0.01\\
\frac{A - A_{\min}}{A_{\max} - A_{\min}} &= \frac{0.0017 - 0.00005}{0.01 - 0.00005} = \frac{0.00165}{0.00995} = 0.1658\\
\mathrm{pAUC}^{\text{McClish}} &= \tfrac12(1 + 0.1658) = 0.5829.
\end{align*}
The same detector reads 0.170 on the repository's scale and 0.583 on McClish's\tagC{}. Never compare the two
without converting.

\begin{worked}[label={we:measures:pauc}]{Partial AUC of the six pixels, with a limit that cuts a segment}
Take the six pixels and a limit $m = 0.5$ (chosen so the cut falls inside the curve; the product uses 0.01). The
segments left of FPR 0.5 are:
\begin{align*}
(0, \tfrac13) \to (\tfrac13, \tfrac13): &\quad \tfrac13 \times \tfrac13 = \tfrac19\\
\text{vertical segments at } \tfrac13: &\quad 0\\
(\tfrac13, 1) \to (\tfrac23, 1), \text{ cut at } 0.5: &\quad (0.5 - \tfrac13) \times 1 = \tfrac16\\
\text{area} &= \tfrac19 + \tfrac16 = \tfrac{2}{18} + \tfrac{3}{18} = \tfrac{5}{18}\\
\mathrm{pAUC}_{0.5} &= \tfrac{5}{18} / 0.5 = \tfrac{10}{18} = \tfrac59 = 0.556.
\end{align*}
Chance at this limit is $0.5 / 2 = 0.25$. And at the product's limit $m = 0.01$: the curve sits at TPR $\frac13$ from
FPR 0 all the way to $\frac13$, so the area up to 0.01 is $0.01 \times \frac13$, and dividing by $0.01$ gives
$\mathrm{pAUC}_{0.01} = \frac13 = 0.333$. The code returns $0.5556$ and $0.3333$\tagC{}.
\end{worked}

\paragraph{Cutting a sloped segment.} When the segment that crosses $m$ is diagonal (a tie, or the straight lines of
detectors A and B), \code{pauc} cuts it by linear interpolation. The segment from $(x_0, y_0)$ to $(x_1, y_1)$ with
$x_0 < m < x_1$ has, at $m$, the height
\begin{equation}\label{eq:measures:interp}
y(m) = y_0 + (y_1 - y_0)\,\frac{m - x_0}{x_1 - x_0},
\end{equation}
and its piece left of $m$ is the trapezoid of width $m - x_0$ and heights $y_0$, $y(m)$. For the tie of
\cref{we:measures:tie} with $m = 0.25$: the diagonal runs $(0, \frac12) \to (\frac12, 1)$, so
$y(0.25) = 0.5 + 0.5 \times \frac{0.25}{0.5} = 0.75$; the area is $0.25 \times \frac{0.5 + 0.75}{2} = 0.15625$; and
$\mathrm{pAUC}_{0.25} = 0.15625 / 0.25 = 0.625$. For detector A at $m = 0.01$ the cut falls exactly at a point, so
$\mathrm{pAUC}_A = 0.002 / 0.01 = 0.20$ and $\mathrm{pAUC}_B = 0.0001 / 0.01 = 0.01$: the order flips to the one the
analyst would choose.

The code, annotated:
\begin{lstlisting}[basicstyle=\ttfamily\footnotesize]
def pauc(fpr, tpr, max_fpr=0.01):            # fpr, tpr: (K,) exact ROC points
    x0, x1, y0, y1 = fpr[:-1], fpr[1:], tpr[:-1], tpr[1:]          # (K-1,) segments
    y1c = np.where(x1 > max_fpr,                                    # segment crosses the limit:
                   y0 + (y1 - y0) * (max_fpr - x0) / np.where(x1 > x0, x1 - x0, 1.0),  # eq. interp
                   y1)
    return float(np.sum(((np.minimum(x1, max_fpr) - x0) * (y0 + y1c) / 2)[x0 < max_fpr]) / max_fpr)
\end{lstlisting}
Segments that start at or beyond the limit (\code{x0 >= max_fpr}) are dropped; vertical runs at one FPR (several
targets above every background pixel) have zero width but do raise the height that later trapezoids use, so they
``count fully'', as the docstring says.

\section{Exact ROC versus the product's sampled ROC}\label{sec:measures:exact}

The product's \code{compute_roc} (\file{app/utils/anomaly_detection/scoring.py}, line 184) does not use every score
as a threshold. It takes evenly spaced percentiles:
\begin{lstlisting}
n_pts = min(max_thresholds, max(20, s.size // 10000 + 64))      # max_thresholds = 256
qs = np.linspace(0.0, 100.0, n_pts)
\end{lstlisting}
For a scene of 1,000,000 valid pixels:
\begin{align*}
n_{\text{pts}} &= \min(256,\ \max(20,\ 1{,}000{,}000 \,/\!/\, 10{,}000 + 64)) = \min(256, \max(20, 100 + 64)) = 164\\
\text{spacing} &= 100 / (164 - 1) = 100 / 163 = 0.613 \text{ percentile.}
\end{align*}
(Here $/\!/$ is integer division.) So between FPR 0 and 1\% the sampled curve has about one point; its area there is
a single triangle guessed from it. For the whole AUC the coarseness matters little; for the first 1\% it is most of
the answer.

\begin{sourcenote}[: the sampled ROC]
Two descriptions disagree with \code{compute_roc}, and the code wins. ``Samples 256 thresholds'' is a cap: the
count reaches 256 only from $(256 - 64) \times 10{,}000 = 1{,}920{,}000$ pixels. The function's own docstring says
the thresholds are ``evenly spaced over the score range plus the unique-value breakpoints sklearn would emit''; the
code takes percentiles only (\file{app/utils/anomaly_detection/scoring.py}, lines 184--185).
\end{sourcenote}

\code{exact_roc} avoids sampling. It sorts every score, highest first, and makes one ROC point at the end of each run
of tied scores:
\begin{lstlisting}[basicstyle=\ttfamily\footnotesize,literate={~}{{\textasciitilde}}1]
def exact_roc(s, g):                         # s: (n,) scores, g: (n,) bool truth
    order = np.argsort(-s, kind="stable")
    s_sorted, g_sorted = s[order], g[order]
    last = np.r_[np.flatnonzero(np.diff(s_sorted)), len(s_sorted) - 1]   # end of each tie run
    fpr = np.r_[0.0, np.cumsum(~g_sorted)[last] / max((~g).sum(), 1)]
    tpr = np.r_[0.0, np.cumsum(g_sorted)[last] / max(g.sum(), 1)]
    return fpr, tpr
\end{lstlisting}
That is the walk of \cref{we:measures:six}, for a million pixels: the cumulative sums are the TP and FP columns.
The evaluator reports the product's sampled AUC as \code{roc_auc} (for continuity with every older table) and uses
the exact curve for everything else; the test checks \code{auc_exact} against scikit-learn's \code{roc_auc_score}
to $10^{-9}$.

\section{Detection rate at 0.1\% false alarms}\label{sec:measures:tpr}

\begin{definition}[TPR at a fixed FPR]
The \index{TPR at FPR}\emph{TPR at FPR $f_0$} is the best detection rate reached at any exact ROC point whose FPR
is at most $f_0$:
\begin{equation}\label{eq:measures:tprat}
\mathrm{TPR}@f_0 = \max\{\mathrm{TPR}_k : \mathrm{FPR}_k \le f_0\}.
\end{equation}
In code, \code{float(tpr[fpr <= max_fpr].max())}; the repository's key is \code{tpr_at_0.001}, $f_0 = 0.001$.
\end{definition}

There is no interpolation: only actual ROC points count. On a million-pixel scene 0.1\% of the background is about
1,000 false alarms, roughly what an analyst can inspect. Its chance level is $f_0$ itself, 0.001, because on the
diagonal TPR equals FPR. For the six pixels, the ROC points with FPR at most 0.001 are $(0,0)$ and $(0, \frac13)$,
so $\mathrm{TPR}@0.001 = \frac13$.

\section{Objects caught at a fixed budget}\label{sec:measures:objects}

Pixel rates have a weakness. A 20-pixel roof and a 1-pixel vehicle are two objects, but in TPR the roof counts
twenty times. An analyst needs \emph{one} flagged pixel per object to go and look.

\begin{definition}[Object]
An \index{object}\emph{object} is a \index{connected component}\emph{connected component} of the ground-truth mask:
a maximal set of target pixels in which you can walk from any pixel to any other by steps up, down, left or right
through target pixels. This is \emph{4-connectivity}: two target pixels that touch only at a corner are two
objects. \code{PrismaBench} labels objects with \code{ndimage.label(gt)}, whose default structure is the
4-neighbour cross (\file{app/evaluation/datasets.py}, line 160); implant benches store their object ids directly.
\end{definition}

\begin{definition}[Objects caught at the top $q$\%]
Flag the valid pixels scoring at or above the $(100 - q)$th percentile of all valid scores. Grow the flagged set by
one 4-neighbour step (a \index{binary dilation}\emph{binary dilation}). An object is \emph{caught} when any of its
pixels lies in the grown set. The measure is the share of objects caught.
\end{definition}

The code is three lines (\file{app/evaluation/metrics.py}, \code{object_hits}):
\begin{lstlisting}[basicstyle=\ttfamily\footnotesize]
k = valid & np.isfinite(score)                                    # (H, W) usable pixels
# flagged pixels, grown by one 4-neighbour step (the 1-pixel tolerance)
hits = ndimage.binary_dilation((score >= np.percentile(score[k], 100 - top_pct)) & k)
return set(np.unique(objects[hits & (objects > 0)]).tolist())    # ids of objects touched
\end{lstlisting}

\begin{sourcenote}[: the rule is ``within one pixel'', not ``a pixel above the threshold'']
You will read that an object is caught when ``at least one of its pixels is above the threshold''. The code is
more generous: a flagged pixel \emph{next to} the object (one 4-neighbour step) also catches it. This is the
\index{1-pixel tolerance}\emph{1-pixel tolerance}; the module docstring of \file{app/evaluation/metrics.py} says
``objects with a pixel within 1 px of the top top\_pct\% of scores''. The reason is that the ground-truth masks may
sit about one pixel off the scene grid. Note also that the percentile is taken over \emph{all} valid pixels, targets
included, so ``top 1\%'' is not exactly ``FPR 1\%''. Pixel measures (AUC, pAUC, TPR) have no tolerance.
\end{sourcenote}

\begin{figure}[htbp]
\centering
\begin{tikzpicture}[x=0.55cm,y=0.55cm]
\foreach \x in {0,...,6} \foreach \y in {0,...,4} \draw[gray!40] (\x,\y) rectangle ++(1,1);
% domino object at (2,2),(3,2)
\fill[spat!25] (2,1) rectangle ++(1,1); \fill[spat!25] (3,1) rectangle ++(1,1);
\fill[spat!25] (2,3) rectangle ++(1,1); \fill[spat!25] (3,3) rectangle ++(1,1);
\fill[spat!25] (1,2) rectangle ++(1,1); \fill[spat!25] (4,2) rectangle ++(1,1);
\fill[loss!70] (2,2) rectangle ++(1,1); \fill[loss!70] (3,2) rectangle ++(1,1);
\node at (3,-0.7) {\small domino: $m = 2 + 6 = 8$};
\begin{scope}[xshift=5cm]
\foreach \x in {0,...,6} \foreach \y in {0,...,4} \draw[gray!40] (\x,\y) rectangle ++(1,1);
\fill[loss!70] (2,2) rectangle ++(1,1); \fill[loss!70] (3,3) rectangle ++(1,1);
\node at (2.5,2.5) {\small 1}; \node at (3.5,3.5) {\small 2};
\node at (3.5,-0.7) {\small corner touch: two objects};
\end{scope}
\end{tikzpicture}
\caption{Left: a two-pixel object (amber) and the six pixels its 1-pixel tolerance adds (indigo); a flagged pixel
anywhere in these $m = 8$ cells catches it. Right: two target pixels touching only at a corner are two objects under
4-connectivity (the book's script confirms \texttt{ndimage.label} returns 2).}
\label{fig:measures:dilate}
\end{figure}

\section{The chance level of objects caught}\label{sec:measures:chance}

Here is the trap the tolerance creates. Give every pixel an independent random score. At the top-$q$\% budget, each
pixel is flagged with probability $p = q/100$. Write $m$ for the object's size \emph{after} the 1-pixel dilation:
its own pixels plus its 4-neighbours. The object is caught if any of those $m$ cells is flagged.

\paragraph{Derivation.} The chance of catching it is one minus the chance of missing all $m$ cells:
\begin{align*}
P(\text{one cell not flagged}) &= 1 - p\\
P(\text{none of $m$ flagged}) &= (1 - p)^m && \text{independent cells: multiply}\\
P(\text{caught}) &= 1 - (1 - p)^m. && \text{the complement}
\end{align*}
This is \code{chance_hit(m, top_pct)} in \file{app/evaluation/metrics.py}, and the evaluator stores it beside every
object's hit (\code{chance_1}, \code{chance_0.1}). It is a slight approximation: a percentile threshold flags
exactly $q$\% of pixels, so the cells are not perfectly independent; with a million pixels the difference is
negligible.

\begin{worked}[label={we:measures:chance}]{Chance for four object shapes at the top 1\%}
With $p = 0.01$:
\begin{align*}
\text{1 pixel: } m &= 1 + 4 = 5, & 0.99^5 &= 0.9510, & 1 - 0.9510 &= 0.0490\\
\text{domino: } m &= 2 + 6 = 8, & 0.99^8 &= 0.9227, & 1 - 0.9227 &= 0.0773\\
\text{3 in a row: } m &= 3 + 8 = 11, & 0.99^{11} &= 0.8953, & 1 - 0.8953 &= 0.1047\\
3 \times 3 \text{ square: } m &= 9 + 12 = 21, & 0.99^{21} &= 0.8097, & 1 - 0.8097 &= 0.1903
\end{align*}
For the square, the 12 neighbours are 3 on each of its 4 sides. At the top 0.1\% ($p = 0.001$) the one-pixel object
has $1 - 0.999^5 = 1 - 0.99501 = 0.0050$. The script's values match: 0.0490, 0.0773, 0.1047, 0.1903 and 0.0050\tagC{}.
\end{worked}

Bigger objects are caught by chance more often, because they have more cells. So the chance level of a whole bench
depends on its size mix. The evaluator averages \code{chance_hit} over the objects and prints it as the
\code{*chance*} row of every table: \textbf{10.2\%} at the top 1\% on the PRISMA bench and \textbf{10.0\%} on
bench-v1\hfil\penalty50\hfilneg\tagM{docs/training/anomaly-detection/19-running-an-evaluation.md}, about 1\% at the top 0.1\%.

\begin{keypoint}
Read every ``objects caught'' number next to its chance row: about 10\% at the top 1\% comes free.
\end{keypoint}

\section{The 3D-ROC terms, and why they are degenerate here}\label{sec:measures:3d}

The hyperspectral literature also reports a ``3D ROC'', which plots TPR and FPR each against the threshold. Scores are
first \emph{min-max normalised}, $\tilde s = (s - s_{\min})/(s_{\max} - s_{\min})$, so the threshold $\tau$ runs from
0 to 1. Two areas follow: $\mathrm{AUC}(D, \tau)$ under detection rate against $\tau$, and $\mathrm{AUC}(F, \tau)$
under false-alarm rate against $\tau$ (lower is better). From them, $\mathrm{AUC}_{\mathrm{BS}} =
\mathrm{AUC} - \mathrm{AUC}(F,\tau)$ (background suppressibility) and $\mathrm{SNPR} = \mathrm{AUC}(D,\tau) /
\mathrm{AUC}(F,\tau)$.

\paragraph{The identity the code uses.} For scores in $[0, 1]$, the area under ``share of pixels with $\tilde s \ge
\tau$'' as $\tau$ runs from 0 to 1 is the \emph{mean} normalised score. Proof: a pixel with normalised score $u$ is at
or above every threshold from 0 to $u$, so it contributes $\int_0^1 \mathbf{1}[u \ge \tau]\,d\tau = u$ to the area;
averaging over pixels gives the mean. So \code{three_d_roc} returns the target mean and the background mean of
$\tilde s$. On the course's five pixels ($0.9$ T, $0.8$ B, $0.7$ T, $0.1$ B, $0.0$ B) the normalised scores are
$1, 0.889, 0.778, 0.111, 0$; so $\mathrm{AUC}(D,\tau) = (1 + 0.778)/2 = 0.889$ and $\mathrm{AUC}(F,\tau) = (0.889 +
0.111 + 0)/3 = 0.333$, as the test asserts.

\paragraph{Why it does not help for residual Mahalanobis.} That score is a squared distance with a heavy tail: a few
pixels sit hundreds of times above the typical background. Min-max normalisation divides by the single highest
score. With a typical background of 30 and one pixel at 30,000\tagA{}, the normalised background is about
$30/30{,}000 = 0.001$ for every model, so $\mathrm{AUC}_{\mathrm{BS}}$ is ROC-AUC minus 0.001 and SNPR is set by the
one brightest pixel. The FT02b README measured exactly that on the PRISMA bench (``AUC\_BS = ROC-AUC - 0.001''). The
terms are computed for comparison with the literature; do not choose a model with them.

\section{The measures side by side}\label{sec:measures:table}

\begin{center}
\small
\begin{tabular}{p{3.0cm}p{2.6cm}ccp{4.1cm}}
\toprule
measure & key & perfect & chance & question it answers \\
\midrule
ROC-AUC & \texttt{roc\_auc} (sampled), exact per group & 1 & 0.5 & are targets above background on average over everything? \\
pAUC to 1\% FPR & \texttt{pauc\_0.01} & 1 & 0.005 & how good is the top of the list? \\
TPR at 0.1\% FPR & \texttt{tpr\_at\_0.001} & 1 & 0.001 & what share of target pixels is in about the top 1,000? \\
caught, top 1\% & \texttt{hit\_1} & 100\% & about 10\% & how many objects in a 1\% review budget? \\
caught, top 0.1\% & \texttt{hit\_0.1} & 100\% & about 1\% & the same, in a 0.1\% budget \\
3D-ROC terms & \texttt{auc\_d\_tau}, \ldots & - & - & literature comparison only \\
\bottomrule
\end{tabular}
\end{center}

\Cref{ch:benchmarks} explains why the four labelled PRISMA scenes could not separate models even with these
measures, and builds a bigger benchmark; \cref{ch:runeval} runs them all from one config; \cref{ch:bootstrap} asks
how large a difference in them is real.

\begin{confusions}
\confusion{A pAUC of 0.15 is terrible - it is far below 0.5?}{The chance level of the repository's pAUC is
$m/2 = 0.005$, not 0.5. A pAUC of 0.15 is thirty times chance. Only McClish's standardisation (not used here) puts
chance at 0.5.}
\confusion{\texttt{roc\_auc} and the exact AUC should be the same number?}{\code{roc_auc} is the product's sampled
\code{compute_roc}, kept so results match what users see; the exact AUC uses every score. They differ slightly;
quote which one you used.}
\confusion{Top 1\% means exactly 1\% of pixels are flagged?}{It means ``at or above the 99th percentile of valid
scores''. With many tied scores more can pass. A score map with a block of zeros (masked pixels left in) can put the
99th percentile at 0, so every pixel is ``flagged'' and every object ``caught'' whatever the scores.}
\confusion{Objects caught at 10\% means the model found one object in ten?}{It does, and so would random scores. Report it
next to chance.}
\confusion{The 1-pixel tolerance applies to AUC too?}{No. Only objects caught has it. AUC, pAUC and TPR at a fixed
FPR are pixel measures with no tolerance.}
\confusion{A diagonal pair of target pixels is one object?}{No. Objects are 4-connected; a corner touch makes two
objects, each with its own chance level.}
\end{confusions}

\begin{checkyourself}
\item\label{ex:measures:1} Four pixels: scores 0.6 (T), 0.5 (B), 0.4 (B), 0.3 (T). List the ROC points, compute the
AUC by trapezoids and by counting pairs, and compute $\mathrm{pAUC}_{0.01}$ and $\mathrm{TPR}@0.001$. What do the
measures disagree about?
\item\label{ex:measures:2} Derive the chance level of the normalised pAUC at a limit $m = 0.05$, and the McClish
value of a raw partial area $A = 0.02$ at that limit.
\item\label{ex:measures:3} With random scores at the top-1\% budget, what is the chance of catching an L-shaped
three-pixel object (cells (0,0), (1,0), (1,1))? Count its 1-pixel tolerance cells by hand.
\item\label{ex:measures:4} A scene has 600,000 valid pixels. How many thresholds does \code{compute_roc} use, and
how many percentiles apart are they? Roughly how many sampled points fall between FPR 0 and 1\%?
\item\label{ex:measures:5} A bench has equal numbers of one-pixel and two-pixel (domino) objects. A model catches
12\% of them at the top 1\%. How many times chance is that?
\item\label{ex:measures:6} Using the staircase proof of \cref{sec:measures:mw}, compute the AUC of
\cref{ex:measures:1} a third way: one rectangle per background pixel.
\end{checkyourself}

\nextchapter{ch:benchmarks}
